\documentclass[portuguese]{sbc2025}%

%\usepackage{graphicx}% já está incluido na classe
%\usepackage[utf8]{inputenc} % obsoleto, desnecessário

\usepackage[misc,geometry]{ifsym} 
\usepackage{hyperref}

%%%%  %\usepackage{fontspec}

%% Os problemas com a classe eram originados por carregarem ambos os
%% pacotes fontenc e fontspec simultaneamente. Agora a classe detecta
%% qual engine está em uso e se for o pdflates então carrega o
%% fontenc. Se forem o xelatex ou lualatex então carrega o
%% fontspec. Eu particularemnte prefiro o lualatex ou o xetex


% \usepackage{fontawesome} %%% incluido na classe. Desnecessário
% chamá-lo aqui


%%%% \usepackage{academicons}
%%%% outro encrenqueiro que se tornou desnecessario. Deste typeface só
%%%% se usava o glifo do orcid e o codigo interno bombava.
%%%% Substitui pelo pacote orcidlink (carregado internamete na classe)
%%%% e consertei o codigo na classe.


%\usepackage{color} % carregada dentro da classe
%\usepackage{hyperref} % carregada dentro da classe

\usepackage{aas_macros}
\usepackage[bottom]{footmisc}
\usepackage{float}
\usepackage{booktabs}

%\usepackage{supertabular}
%\usepackage{multicol}
%\usepackage{multirow}
%% O pacote tabularray é muito melhor para
% a criacao de tabela complexas e/ou loooongas; tudo de modo simples.
%% Torna o uso de multicol e multirow desnecessarios

\usepackage{tabularray}


\usepackage{afterpage}
\usepackage{url}
\usepackage{pifont}


\setcitestyle{square}


\definecolor{engtitle}{rgb}{0.5,0.5,0.5}
\definecolor{orcidlogo}{rgb}{0.37,0.48,0.13}
\definecolor{unilogo}{rgb}{0.16, 0.26, 0.58}
\definecolor{maillogo}{rgb}{0.58, 0.16, 0.26}
\definecolor{darkblue}{rgb}{0.0,0.0,0.0}
\hypersetup{colorlinks,breaklinks,
            linkcolor=darkblue,urlcolor=darkblue,
            anchorcolor=darkblue,citecolor=darkblue}
%\hypersetup{colorlinks,citecolor=blue,linkcolor=blue,urlcolor=blue}

%%%%%%% IMPORTANT: We disable hyperlinks by default with this line, to avoid the error "\pdfendlink ended up in different nesting level" while writing.
%\hypersetup{draft}

\category{Artigo de Pesquisa/Research Paper}
\title[Template SOL Preprints para trabalhos em português]{Otimização de intervenções viárias na malha de Brasília com base no índice de desempenho de velocidade}
%\engtitle{\textcolor{engtitle}{Aprimorando o planejamento viário com grafos e dados reais}}

%THE ORCID IS MANDATORY FOR EACH AUTHOR IN JBCS
\author[Marcinoni et al. 2026]{

\affil{\textbf{Pedro Marcinoni Varão Ribeiro Vieira da Rosa}
\\Universidade de Brasília~\ror{02xfp8v59}
\\\href{mailto:pedroextrer@gmail.com}{\textit{pedroextrer@gmail.com}}}

\affil{\textbf{Davi Brasileiro Gomes}
\\Universidade de Brasília~\ror{02xfp8v59}
\\\href{mailto:davi.brasileiro.gomes@gmail.com}{\textit{davi.brasileiro.gomes@gmail.com}}}

\affil{\textbf{Davi Galvão Guerra}
\\Universidade de Brasília~\ror{02xfp8v59}
\\\href{mailto:daviguerra27@gmail.com}{\textit{daviguerra27@gmail.com}}}

\affil{\textbf{Leonardo Mileo Moreira Krauss Guimarães}
\\Universidade de Brasília~\ror{02xfp8v59}
\\\href{mailto:leonardokrauss3107@gmail.com}{\textit{leonardokrauss3107@gmail.com}}}

\affil{\textbf{Patrícia Araújo de Oliveira}
\\Universidade de Brasília~\ror{02xfp8v59}
\\\href{patricia.oliveira@unb.br}{\textit{patricia.oliveira@unb.br}}}

}


\begin{document}

\begin{frontmatter}

\maketitle

\begin{abstract-pt}
Este artigo investiga o impacto de intervenções viárias sobre o congestionamento e os tempos de deslocamento, além do uso de grafos direcionados para identificar gargalos na malha urbana. Para isso, modelou-se a região do trevo EPCT/EPIP/DF-480, utilizando dados do OpenStreetMap e da TomTom. O tráfego foi simulado por meio do algoritmo de Frank-Wolfe, considerando o equilíbrio de Wardrop e funções volume-atraso BPR. Foram avaliados cenários de adição, remoção e inversão de vias, comparando seus efeitos sobre o Tempo Total de Viagem do Sistema e o Índice de desempenho de velocidade. [RESULTADOS: inserir os principais resultados quantitativos das intervenções analisadas.] [CONCLUSÃO: sintetizar o que os resultados demonstram sobre os efeitos das intervenções e a utilidade da modelagem por grafos.]
\end{abstract-pt}

\begin{keywords}
Índice de desempenho de velocidade, Tempo Total de Viagem do Sistema, Equilíbrio de Wardrop, Função volume-atraso, Paradoxo de Braess
\end{keywords}


%\begin{license}
%Published under the Creative Commons Attribution 4.0 International Public License (CC BY 4.0)
%\end{license}

\end{frontmatter}


\section{Introdução}
\label{sec:intro}
O crescimento da demanda por mobilidade urbana exige um planejamento viário capaz de reduzir congestionamentos e melhorar os tempos de deslocamento. Nesse contexto, a identificação de gargalos e a avaliação dos efeitos de intervenções na malha viária são fundamentais para compreender o desempenho do sistema. Este estudo concentra-se na região do Gama, no Distrito Federal, especialmente no entroncamento entre a EPCT, a EPIP e a DF-480.

A pesquisa busca responder a duas questões: \textit{RQ1 — Como intervenções pontuais na malha viária afetam os tempos de deslocamento e o congestionamento nas vias adjacentes? RQ2 — Como a modelagem da infraestrutura viária por meio de grafos direcionados ponderados contribui para identificar gargalos e compreender as causas dos congestionamentos?}

Para respondê-las, são integrados dados geográficos e de tráfego em um modelo de rede viária, permitindo simular intervenções de adição de faixas, remoção de ligações e inversão de sentido. Seus efeitos são avaliados por meio de indicadores de fluidez e tempo total de viagem, considerando tanto os trechos modificados quanto o desempenho global da rede e possíveis diferenças entre esses indicadores. O artigo, portanto, é dividido nos seguintes tópicos: II. Fundamentação Teórica, III. Metodologia, IV. Resultados e V. Conclusão.
 
\section{Fundamentação Teórica}
Aqui são explorados conceitos técnicos associados à análise de congestionamento viário, tais como o Paradoxo de Braess, Equilíbrio de Wardrop, Função Volume-Atraso e Indicador de Performance de Velocidade.

\subsection{Índice de desempenho de velocidade (K/SPI)}
O Speed Performance Indicator (SPI) é uma métrica que analisa o desempenho viário pela razão entre a velocidade média observada em uma via, e a velocidade máxima permitida nela \cite{he2016spi}. Aqui, o chamamos de K por fins de facilidade.

\subsection{Equilíbrio do usuário de Wardrop}
O equilíbrio de Wardrop é definido em uma situação onde nenhum motorista pode, individualmente, alterar seu tempo médio de deslocamento na malha ao mudar o trajeto escolhido \cite{wardrop1952}.

\subsection{Função volume-atraso (BPR)}
A Função Volume-Atraso (VDF) representa a variação do tempo de deslocamento considerando a razão entre volume de veículos e a capacidade viária na malha observada. Matematicamente, é modelada como $t(v) = t_0 . (1 + a(v/c)^b)$, sendo os parâmetros $a$ e $b$ coeficientes de calibração específicos de cada malha.

\subsection{Paradoxo de Braess}
Define uma situação na qual adicionar uma nova via em uma malha pode não reduzir o congestionamento \cite{braess2005}, mas sim aumentá-lo e criar novos gargalos.

\section{Metodologia}
\label{sec:metodologia}

Esta seção apresenta uma descrição sobre a modelagem em grafo da área de estudos selecionada, a obtenção dos dados relacionados ao sistema viário da região, a implementação do algoritmo de alocação de tráfego e a formulação das propostas de intervenções no grafo.

\subsection{Visão geral e área de estudo}

Adotou-se um modelo de alocação estática de tráfego, representando a malha viária por um grafo direcionado com demanda fixa. As intervenções foram realizadas por alterações discretas no grafo \cite{leblanc1975dndp}, enquanto a redistribuição dos fluxos foi simulada pelo algoritmo de Frank-Wolfe \cite{leblanc1975fw}, considerando o equilíbrio de Wardrop \cite{wardrop1952} e a função volume-atraso BPR \cite{bpr1964}.

A área de estudo compreende o entorno do trevo entre a EPCT (DF-001), a EPIP e a DF-480, no Gama (DF). A rede foi obtida do OpenStreetMap por meio do OSMnx 2.1.1. O subgrafo inicial de 30 nós foi ampliado para incorporar rotas alternativas, utilizando vias de classe terciária ou superior contidas na envoltória convexa dos nós originais, com margem de 200~m. O grafo resultante possui 134 nós e 236 arestas, com 40,2\% da demanda dispondo de pelo menos duas rotas competitivas.

\subsection{Demanda}

A demanda foi obtida pela ferramenta TomTom O/D Analysis, com regiões de origem e destino associadas aos 134 nós da rede e registros de 22 dias úteis de julho de 2026, entre 18h e 19h. As viagens externas foram associadas à primeira ou à última região atravessada. A partir dos percentuais fornecidos pela plataforma, foram reconstruídas 15.429 viagens, das quais 86\% correspondiam a tráfego de passagem. Após a exclusão de viagens intranó e pares sem caminho, restaram 14.960 viagens em 1.031 pares origem-destino.

Como a amostra representa apenas parte do tráfego real, aplicou-se um fator de expansão calibrado com a mediana das contagens volumétricas do DER-DF para o mesmo período e horário. Calculado pela razão entre o volume observado e o fluxo da amostra em quatro pontos de contagem, o fator variou de 8,82 a 21,96. Adotou-se a mediana de 12,21 como cenário-base (8.309 veículos/h), utilizando os extremos para representar os níveis de demanda baixa e alta.

\subsection{Atributos das arestas}

As velocidades foram obtidas do TomTom Traffic Stats para os dias úteis de julho de 2026, nos períodos de pico (18h--19h) e madrugada (0h--4h). Os registros foram associados às arestas conforme sua posição e sentido, calculando-se a velocidade pela média harmônica ponderada pelo comprimento sobreposto. O índice $K$ foi definido pela razão entre a velocidade de pico e o limite regulamentado da via \cite{he2016spi}.

Os limites de velocidade e o número de faixas foram extraídos do OSM e, quando ausentes, estimados com base nas características das vias e no Código de Trânsito Brasileiro \cite{ctb2026}. A capacidade de cada aresta foi calculada pelo número de faixas multiplicado pela capacidade por faixa, entre 600 e 1.800 veículos/h, conforme a classe viária. O tempo em fluxo livre foi estimado pela maior velocidade entre os períodos de pico e madrugada, limitada ao valor regulamentado.

Como a função BPR \cite{bpr1964}, com parâmetros clássicos ($a=0{,}15$ e $b=4$), subestimou os atrasos observados no trevo, foram ajustadas as capacidades efetivas de 25 arestas em que o fluxo estimado superava 30\% da capacidade nominal e o atraso observado excedia o previsto. O ajuste buscou reproduzir os tempos observados, respeitando o limite da capacidade nominal, embora parte da avaliação de consistência permaneça dependente da calibração.

\subsection{Alocação, cenários e indicadores}

O equilíbrio de Wardrop foi obtido pela formulação de Beckmann \cite{beckmann1956}, sujeita às restrições de conservação de fluxo:
\begin{equation}
\min_{\mathbf{x}}\; \sum_{e}\int_{0}^{x_e} t_e(w)\,dw,\quad t_e(x)=t_{0,e}\Big[1+a\Big(\tfrac{x}{c_e}\Big)^{b}\Big].
\label{eq:beckmann}
\end{equation}
O algoritmo de Frank-Wolfe \cite{frank1956} parte da alocação em fluxo livre e atualiza iterativamente os tempos e caminhos mínimos, redistribuindo os fluxos até que o \textit{gap} relativo, que mede a diferença entre o custo da alocação atual e o dos caminhos mínimos sob os tempos vigentes, seja inferior a $10^{-4}$ ou o limite de 500 iterações seja atingido. O desempenho foi avaliado pelo índice $K_e=v_e/v_{\text{lim},e}$, com $v_e=L_e/t_e$, e pelo Tempo Total de Viagem do Sistema ($TSTT=\sum_e x_e t_e(x_e)$).

Foram simuladas intervenções de adição de faixas, remoção de arestas e inversão de sentido em trechos selecionados a partir dos indicadores de fluidez e fluxo. Seus efeitos foram comparados ao cenário-base, investigando-se possíveis deteriorações sistêmicas associadas ao Paradoxo de Braess \cite{braess2005}. As simulações consideraram três níveis de demanda e uma análise de sensibilidade aos parâmetros BPR ($a=0{,}8$ e $b=4{,}7$) \cite{vdf2023}, com as capacidades efetivas recalculadas.

\subsection{Limitações do desenho}

O modelo estático não considera a propagação de filas, os atrasos em semáforos ou rotas por vias locais, excluídas do grafo. Além disso, a predominância de viagens de passagem limita as alternativas de rota disponíveis. As capacidades viárias foram estimadas com base na literatura, enquanto parte dos limites de velocidade e do número de faixas foi preenchida por estimativas. Por fim, a avaliação de consistência não constitui uma validação independente, pois os dados observados também foram utilizados na construção e calibração do modelo.

\section{Resultados}

\subsection{Diagnóstico inicial do K}
\label{sec:diagnostico_k}

A Fig.~\ref{fig:k_inicial} apresenta a distribuição espacial do índice $K$ no cenário-base, considerando o horário de pico (18h--19h).

\begin{figure}[htbp]
\centering
\includegraphics[width=6cm]{gama_drive_expanded_k_inicial.png}
\caption{Distribuição espacial do índice $K$ na região do Gama (EPCT/EPIP/DF-480).}
\label{fig:k_inicial}
\end{figure}

A análise revela diferentes condições de circulação: as rodovias principais apresentam valores de $K$ superiores a $0{,}8$, indicando velocidades próximas aos limites regulamentados; as vias coletoras e conexões secundárias apresentam valores entre $0{,}4$ e $0{,}7$, indicando redução moderada das velocidades; e os gargalos mais críticos, que apresentam valores de $K$ inferiores $0{,}3$, a concentram-se no trevo entre a EPCT, a EPIP e a DF-480, especialmente nas alças de acesso.

\subsection{Convergência do Algoritmo e Avaliação de Consistência}
O algoritmo de Frank-Wolfe foi empregado para obter uma solução aproximada do Equilíbrio de Wardrop na rede viária modelada. A Tabela~\ref{tab:convergencia_base} sintetiza a convergência computacional obtida nos três níveis de demanda no cenário-base.

\begin{table}[htbp]
\caption{Desempenho de Convergência do Algoritmo de Frank-Wolfe}
\label{tab:convergencia_base}
\centering
\resizebox{\columnwidth}{!}{%
\begin{tabular}{lcccc}
\hline
\textbf{Nível de Demanda} & \textbf{Fator} & \textbf{Iterações} & \textbf{\textit{Gap} Final} & \textbf{Tempo (s)} \\
\hline
Baixa & 8,82 & 2 & $3,5 \times 10^{-16}$ & 0,21 \\
Adotada & 12,21 & 4 & $3,1 \times 10^{-9}$ & 0,34 \\
Alta & 21,96 & 28 & $1,0 \times 10^{-4}$ & 2,23 \\
\hline
\end{tabular}%
}
\end{table}

A conferência com os dados reais de tráfego (nível de demanda adotada) demonstrou forte aderência: o fluxo modelado e as passagens TomTom em 40 arestas isoladas atingiu uma correlação de $0,93$, com $GEH$ mediano de $3,48$, e nas 130 arestas não calibradas, a razão mediana entre tempo modelado e observado foi de $0,912$.

\subsection{Desempenho do Sistema e Efeito de Intervenções}
A Tabela~\ref{tab:indicadores_base} apresenta os principais indicadores do sistema na situação atual da malha viária.

\begin{table}[htbp]
\caption{Indicadores Macro do Sistema}
\label{tab:indicadores_base}
\centering
\resizebox{\columnwidth}{!}{%
\begin{tabular}{lccc}
\hline
\textbf{Indicador} & \textbf{D. Baixa} & \textbf{D. Adotada} & \textbf{D. Alta} \\
\hline
Demanda Total (veículos/h) & 6.002,2 & 8.309,2 & 14.944,3 \\
TSTT (veículos$\cdot$h) & 343,36 & 476,94 & 884,36 \\
Velocidade Média (km/h) & 45,20 & 45,05 & 43,68 \\
Tempo Médio de Viagem (min) & 3,43 & 3,44 & 3,55 \\
Índice de Fluidez Média ($K_{rede}$) & 0,7026 & 0,7001 & 0,6763 \\
Arestas Lentas ($K < 0,5$) & 67 & 67 & 68 \\
\hline
\end{tabular}%
}
\end{table}

Para avaliar o comportamento da malha frente a intervenções de capacidade, simulou-se a adição de uma faixa em 5 trechos de gargalo na EPCT sentido Norte. As alterações relativas ao cenário-base no TSTT, no índice $K$ e na quantidade de arestas saturadas são consolidadas na Tabela~\ref{tab:intervencao_faixa}.

\begin{table}[htbp]
\caption{Impacto da Adição de 1 Faixa na EPCT Norte}
\label{tab:intervencao_faixa}
\centering
\resizebox{\columnwidth}{!}{%
\begin{tabular}{lcccc}
\hline
\textbf{Demanda} & \textbf{$\Delta$ TSTT (\%)} & \textbf{$K$ Base} & \textbf{$K$ Cenário} & \textbf{Arestas $V/C>1$} \\
\hline
Baixa   & -8,59\% & 0,7026 & 0,7266 & 15 (-5) \\
Adotada & -8,57\% & 0,7001 & 0,7241 & 15 (-5) \\
Alta    & -8,32\% & 0,6763 & 0,7003 & 21 (-5) \\
\hline
\end{tabular}%
}
\end{table}

O ganho de eficiência no TSTT situou-se de forma consistente na faixa de $8,3\%$ a $8,6\%$ em todas as hipóteses de demanda. Paralelamente, observou-se um aumento de aproximadamente $0,024$ no $K$ nos três cenários analisados, indicando uma melhoria nas condições gerais de circulação. Além disso, a adição da faixa resultou na redução de cinco arestas saturadas ($V/C > 1,0$) em cada cenário de demanda, evidenciando uma diminuição da quantidade de trechos com fluxo superior à capacidade efetiva considerada no modelo.

\subsection{Redistribuição de Fluxos e Avaliação de Efeitos Sistêmicos}
A ampliação da capacidade nos trechos gargalo da EPCT Norte resultou em aumentos expressivos do índice de fluidez local: na aresta $133 \rightarrow 63$, o índice $K$ subiu de $0,29$ para $0,62$, e na aresta $132 \rightarrow 133$, o acréscimo foi de $+0,221$.

A redução dos custos de deslocamento nos trechos ampliados favoreceu uma redistribuição dos fluxos no novo equilíbrio, com aumento da utilização de determinados acessos à EPCT Norte e redução do fluxo atribuído a rotas alternativas. Na demanda adotada, observou-se um acréscimo de fluxo de até $+6,1$ veículos/h no acesso direto à EFPCT Norte e uma redução concomitante de fluxo em rotas secundárias paralelas.

Não foi identificado um efeito paradoxal associado ao Paradoxo de Braess na intervenção analisada, uma vez que a ampliação de capacidade resultou em redução do tempo total de viagem do sistema. A melhoria local refletiu-se diretamente em uma redução global do tempo do sistema ($\Delta TSTT < 0$) e no aumento do índice de fluidez da rede, confirmando que a expansão de infraestrutura nos gargalos críticos do trevo EPCT/EPIP/DF-480 gera benefícios sistêmicos.
%\subsection{Efeito das intervenções no K e no TSTT}
%\textit{[Comparação entre cenários e linha de base.]}

%\subsection{Redistribuição de fluxo e efeito Braess}
%\textit{[Vias que ganham ou perdem fluxo e casos em que a melhoria local piora o sistema.]}



\section{Discussão e Limitações}
\textit{[Interpretação dos resultados e limitações: escopo rodoviário, matriz OD sintética, cobertura do OSM para faixas, ausência de cortes de rota por vias locais, sensibilidade à capacidade.]}

\section{Conclusão e Trabalhos Futuros}
\textit{[Resposta à RQ em um parágrafo e extensões: otimização com NSGA-II, expansão da região e dados de maior granularidade.]}

\begin{declarations}

\begin{contributions}
PM contribuiu para a definição das etapas de pesquisa, e para a escolha do algoritmo de grafos a ser utilizado nas análises. DB contribuiu para a definição do escopo da análise, e para a criação dos scripts de operacionalização dos dados obtidos. DG e LM conduziram a extração e análise dos dados propriamente dita, sendo responsáveis pela validação das saídas e formatação dos resultados. Todos os autores leram e aprovaram o manuscrito final do artigo.
\end{contributions}

\begin{interests}
Os autores declaram que não houve interesses conflitantes na elaboração do presente artigo, tendo o objetivo principal sido alinhado em torno da submissão ao WPOS considerando um problema real observado por cada integrante.
\end{interests}

\begin{materials}
Os dados e informações obtidas/analisadas no artigo foram obtidos de um híbrido de fontes públicas e privadas. Na primeira categoria, os dados geográficos foram mapeados com base no dataset de \textbf{\href{https://www.dados.df.gov.br/pt/inicio}{dados viários do Departamento de Estradas e Rodagens}}, bem como no dataset de limites geográficos \textbf{\href{https://geo-df-der.hub.arcgis.com/pages/downloads}{dataset de limites geográficos do ArcGis}}. Vale ressaltar que esses dados foram alimentados à biblioteca \textbf{\href{https://osmnx.readthedocs.io/en/stable/}{OpenStreetMapx 2.1.1}}, usada para obtenção de arquivos enriquecidos manipuláveis para o processamento. Quanto à segunda categoria, destaca-se o uso da \textbf{\href{https://www.tomtom.com/products/traffic-apis/}{API de fluxo viário, e de obtenção da Matriz O/D do TomTom}}, que cobra uma taxa mensal, mas oferece um período de teste gratuito, durante o qual foi feita a análise evidenciada. Além disso, fica registrado o uso do \textbf{\href{https://www.python.org}{Python 3.13}} para os scripts de correlação dos dados, de criação dos grafos e de execução do algoritmo de Frank-Wolfe.
\end{materials}

\begin{furtherinformation}
Fica registrado o uso de ferramentas de Inteligência Artificial como suporte para a produção do artigo, especialmente no que tange à seleção, limpeza e formatação dos dados utilizados para a análise. Não obstante, cada etapa da análise teve aprovação manual, e as simulações de intervenção viária foram confirmadas pelos integrantes do grupo.

\end{furtherinformation}

\end{declarations}

\bibliographystyle{apalike-sol}
\begin{thebibliography}{99}
\bibitem{wardrop1952} J. G. Wardrop, ``Some theoretical aspects of road traffic research,'' \textit{Proc. Institution of Civil Engineers}, Part II, vol. 1, no. 3, pp. 325--378, 1952.
\bibitem{beckmann1956} M. J. Beckmann, C. B. McGuire, and C. B. Winsten, \textit{Studies in the Economics of Transportation}. New Haven, CT: Yale University Press, 1956.
\bibitem{frank1956} M. Frank and P. Wolfe, ``An algorithm for quadratic programming,'' \textit{Naval Research Logistics Quarterly}, vol. 3, no. 1--2, pp. 95--110, 1956.
\bibitem{leblanc1975fw} L. J. LeBlanc, E. K. Morlok, and W. P. Pierskalla, ``An efficient approach to solving the road network equilibrium traffic assignment problem,'' \textit{Transportation Research}, vol. 9, no. 5, pp. 309--318, 1975.
\bibitem{leblanc1975dndp} L. J. LeBlanc, ``An algorithm for the discrete network design problem,'' \textit{Transportation Science}, vol. 9, no. 3, pp. 183--199, 1975.
\bibitem{bpr1964} U.S. Bureau of Public Roads, \textit{Traffic Assignment Manual}. Washington, DC: U.S. Dept. of Commerce, 1964.
\bibitem{braess2005} D. Braess, A. Nagurney, and T. Wakolbinger, ``On a paradox of traffic planning,'' \textit{Transportation Science}, vol. 39, no. 4, pp. 446--450, 2005.
\bibitem{he2016spi} F. He, X. Yan, Y. Liu, and L. Ma, ``A traffic congestion assessment method for urban road networks based on speed performance index,'' \textit{Procedia Engineering}, vol. 137, pp. 425--433, 2016.
\bibitem{vdf2023} {Y. Pan, Modified volume-delay function based on traffic fundamental diagram}, \textit{J. Transportation Engineering, Part A}, vol. 149, no. 11, 2023.
\bibitem[Brasil, 1997]{ctb2026} BRASIL. Lei nº 9.503, de 23 de setembro de 1997. Institui o Código de Trânsito Brasileiro. \textit{Diário Oficial da União},Brasília,DF,1997.Disponível em: <http://www.planalto.gov.br/ccivil\_03/leis/L9503.htm>.
\end{thebibliography}

\end{document}
